%%% required packages
\usepackage[utf8]{inputenc}
\usepackage[T1]{fontenc}
\usepackage{xcolor}
\usepackage{booktabs}
\usepackage{babel}
\usepackage{etoolbox}
\usepackage{graphicx}
\usepackage{csquotes}
\usepackage{pifont}
\usepackage[scaled=0.89]{sourcecodepro}
\usepackage{listings}
\usepackage{tikz}
\usepackage{tikz-qtree}
\usepackage{textcomp}
\usepackage{amssymb}
\usepackage{fontawesome}

\usepackage{microtype}
\DisableLigatures[f]{encoding = *, family = tt }

\newcommand{\chck}{\ensuremath\checkmark}
\newcommand{\semester}{Summer 2020}
\newcommand{\veranstaltung}{Title of the Lecture}


\usepackage[shellescape]{gmp}

%%%%%%%
\definecolor{lightyellow}{rgb}{1,1,0.8}
\definecolor{colKeys}{rgb}{0,0,1}
\definecolor{colIdentifier}{rgb}{0,0,0}
\definecolor{colComments}{rgb}{1,0,0}
\definecolor{colString}{rgb}{0,0.5,0}

\definecolor{lightgray}{rgb}{.9,.9,.9}
\definecolor{darkgray}{rgb}{.4,.4,.4}
\definecolor{purple}{rgb}{0.65, 0.12, 0.82}

\definecolor{folderbg}{RGB}{124,166,198}
\definecolor{folderborder}{RGB}{110,144,169}


\setbeamertemplate{enumerate item}{(\alph{enumi})}
\setbeamertemplate{enumerate subitem}{(\roman{enumii})}

% \setbeamertemplate{navigation symbols}{\insertslidenavigationsymbol \insertframenavigationsymbol \insertbackfindforwardnavigationsymbol}

\setbeamertemplate{navigation symbols}{}


\usepackage{forest} % https://tex.stackexchange.com/questions/5073/making-a-simple-directory-tree


\def\Size{4pt}
\tikzset{
  folder/.pic={
    \filldraw[draw=folderborder,top color=folderbg!50,bottom color=folderbg]
      (-1.05*\Size,0.2\Size+5pt) rectangle ++(.75*\Size,-0.2\Size-5pt);  
    \filldraw[draw=folderborder,top color=folderbg!50,bottom color=folderbg]
      (-1.15*\Size,-\Size) rectangle (1.15*\Size,\Size);
  }
}



% https://tex.stackexchange.com/questions/161549/scaling-inline-code-to-the-current-font-size
\makeatletter
\lstdefinestyle{mystyle}{
  basicstyle=%
    \ttfamily
    \color{blue}%
    \lst@ifdisplaystyle\scriptsize\fi
}
\makeatother



\usetikzlibrary{shapes,snakes}

% Venn Diagrams
% Definition of circles
\def\firstcircle{(0,0) circle (2cm)}
\def\secondcircle{(0:3cm) circle (2cm)}

\colorlet{circle edge}{blue!50}
\colorlet{circle area}{blue!20}

\tikzset{filled/.style={fill=circle area, draw=circle edge, thick},
    outline/.style={draw=circle edge, thick}}

%%% end venn

\makeatletter
\lstdefinestyle{ausgabe}{
  basicstyle=\ttfamily\scriptsize,%
  backgroundcolor=\color{lightgray}%
}
\makeatother

\lstnewenvironment{ausgabe}{\lstset{style=ausgabe}}{} 

\lstset{%
    float=hbp,%
    basicstyle=\ttfamily\footnotesize, %
	style=mystyle,%
    identifierstyle=\color{colIdentifier}, %
    keywordstyle=\color{colKeys}, %
    stringstyle=\color{colString}, %
    commentstyle=\color{colComments}, %
    literate={fl}{{f{}l}}2,%
    columns=flexible, %
    tabsize=2, %
    frame=single, %
    upquote=true,%
    extendedchars=true, %
    showspaces=false, %
    showstringspaces=false, %
    numbers=left, %
    numberstyle=\tiny, %
    breaklines=true, %
    backgroundcolor=\color{lightyellow}, %
    breakautoindent=true, %
    captionpos=b%
}

%%%%%%%%%%%%
\lstset{literate=%
    {Ö}{{\"O}}1
    {Ä}{{\"A}}1
    {Ü}{{\"U}}1
    {ß}{{\ss}}1
    {ü}{{\"u}}1
    {ä}{{\"a}}1
    {ö}{{\"o}}1
    {~}{{\textasciitilde}}1
}

\usepackage{MnSymbol} % for $\filledsquare$
\usepackage{svn-multi}
\usepackage[absolute,verbose,overlay]{textpos}

%%%% Color Definitions
\definecolor{fom}{RGB}{0,153,139}
\definecolor{myblue}{RGB}{0,0,125}

\usepackage{mdframed}
\newmdenv [linecolor=myblue, frametitle=Infobox,leftmargin=1cm,rightmargin=1cm]{infobox}

%%% Commands to position the FOM Logo in the top right corner and the horizontal rule in the footer

\newcommand{\poslogo}{%
  \setlength{\TPHorizModule}{1cm}
  \setlength{\TPVertModule}{1cm}
  \begin{textblock}{1}(11.5,0.3)
   \includegraphics[width=1cm]{Template/logo}
  \end{textblock}}
  
  \newcommand{\footaline}{%
  \setlength{\TPHorizModule}{0.5cm}
  \setlength{\TPVertModule}{1cm}
  \begin{textblock}{1}(0.5,9.2)
   {\color{myblue}\rule{\dimexpr\paperwidth-0.6cm\relax}{0.3pt}}
  \end{textblock}}
  
%%% Custom titlepage 

\makeatletter
 \usebackgroundtemplate{
    \ifnum\c@framenumber=1%
        \includegraphics[width=\paperwidth,height=\paperheight]{Template/background}
	\else 
        %\includegraphics[width=\paperwidth,height=\paperheight]{background2}
	 \fi%
}
\makeatother
  
\defbeamertemplate*{title page}{customized}[1][]
{\vspace*{6.5cm}
 \begin{flushright}
  \usebeamerfont{title}\inserttitle\par
  \usebeamerfont{subtitle}\usebeamercolor[fg]{subtitle}\insertsubtitle\par
  \bigskip
  \usebeamerfont{author}\insertauthor\par
  \usebeamerfont{institute}\insertinstitute\par
  \usebeamerfont{date}\insertdate\par
  \usebeamercolor[fg]{titlegraphic}\inserttitlegraphic
\end{flushright}
}
  
%%% Adjustment of the Frame header (color)

\makeatletter
\setbeamertemplate{frametitle}
{
  \ifbeamercolorempty[bg]{frametitle}{}{\nointerlineskip}%
  \@tempdima=\textwidth%
  \advance\@tempdima by\beamer@leftmargin%
  \advance\@tempdima by\beamer@rightmargin%
  \begin{beamercolorbox}[sep=0.3cm,left,wd=\the\@tempdima]{frametitle}
    \vbox{}\vskip-1ex%
    \if@tempswa\else\csname beamer@fteleft\endcsname\fi%
    {\strut\color{myblue}\bfseries\insertframetitle\strut\par%
    {%
      \ifx\insertframesubtitle\@empty%
      \else%
      {\usebeamerfont{framesubtitle}\usebeamercolor[fg]{framesubtitle}\insertframesubtitle\strut\par}%
      \fi
    }%
    \vskip-1.5ex%
    \rule{\dimexpr\paperwidth-1.6cm\relax}{0.4pt}}
    \if@tempswa\else\vskip-.3cm\fi% set inside beamercolorbox... evil here...
  \end{beamercolorbox}%
  \poslogo
    %
}
\makeatother

%%%%%%%%%%%% adjust footnote alignment

\makeatletter
\renewcommand<>\beamer@framefootnotetext[1]{%
  \global\setbox\beamer@footins\vbox{%
    \hsize\framewidth
    \textwidth\hsize
    \columnwidth\hsize
    \unvbox\beamer@footins
    \reset@font\footnotesize
    \@parboxrestore
    \protected@edef\@currentlabel
         {\csname p@footnote\endcsname\@thefnmark}%
    \color@begingroup
      \uncover#2{\@makefntext{%
        \rule\z@\footnotesep\ignorespaces\parbox[t]{.9\textwidth}{#1\@finalstrut\strutbox}\vskip1sp}}%
    \color@endgroup}%
}
\makeatother

%%%%%%%% end adjustment frame header

%%%%%%%% adjustment of footer

\makeatletter
\patchcmd{\beamer@calculateheadfoot}{\advance\footheight by 4pt}{\advance\footheight by 20pt}{}{}
\makeatother

%%%%%%%% end adjustment of footer 


% Color definitions
\setbeamercolor{itemize item}{fg=myblue}
\setbeamercolor{itemize subitem}{fg=myblue}
\setbeamercolor{itemize subsubitem}{fg=myblue}
\setbeamercolor{caption name}{fg=myblue}
\setbeamercolor{description item}{fg=myblue}
\setbeamercolor{enumerate item}{fg=myblue}


%%% setting various things
\setbeamerfont{title}{series=\bfseries,size={\fontsize{18}{20}}}
\setbeamerfont{subtitle}{series=\bfseries}

\setbeamercolor{normal text}{fg=black}
\setbeamercolor{alerted text}{fg=myblue}

\setbeamercolor{frametitle}{fg=myblue}
\setbeamercolor{title}{fg=myblue}
\setbeamercolor{framesubtitle}{fg=myblue}
\setbeamercolor{author in head/foot}{fg=myblue}
\setbeamercolor{section in toc}{fg=myblue}
\setbeamercolor{subsection in toc}{fg=myblue}

\setbeamertemplate{frametitle continuation}{} %\insertcontinuationcount
\setbeamertemplate{section in head/foot}{\hfill\insertsectionheadnumber.~\insertsectionhead}
\setbeamertemplate{section in head/foot shaded}{\color{structure!50}\hfill\insertsectionheadnumber.~\insertsectionhead}
\setbeamertemplate{section in toc}{\inserttocsectionnumber.~\inserttocsection}
\setbeamertemplate{subsection in toc}{\hspace*{1em}\inserttocsectionnumber.\inserttocsubsectionnumber.~\inserttocsubsection\\}


% http://tex.stackexchange.com/questions/152565/define-a-new-block-environment-in-latex-beamer

\newenvironment<>{remark}[1]{%
  \setbeamercolor{block title}{fg=white,bg=myblue}%
  \begin{block}#2{#1}}{\end{block}}
\newenvironment<>{examplesecond}[1]{%
  \setbeamercolor{block title}{fg=white,bg=blue!75!black}%
  \begin{block}#2{#1}}{\end{block}}

\setbeamertemplate{navigation symbols}{} % no navigation symbols 



\setbeamertemplate{itemize item}{$\filledsquare$} % from MnSymbols
\setbeamertemplate{itemize subitem}{\textbullet} % if you wnat a circle
\setbeamertemplate{itemize subsubitem}{\textendash}% if you want a triangle


% TOC Definition

\AtBeginSection[]
{
 \begin{frame}<beamer>
 \frametitle{Overview of this Lecture}
 \framesubtitle{~}
 
\tableofcontents[%
    currentsection,% 
    sectionstyle=show/hide,% 
	hideothersubsections
    ]%
    \addtocounter{framenumber}{-1}
  \end{frame}
}

%%%%%%%%%%%%%%%%%%%%%%%%%%%%%

\makeatother
\setbeamertemplate{footline}
{%
  \leavevmode%
%   \rule{\dimexpr\paperwidth-0.6cm\relax}{0.4pt}}
   \hbox{\begin{beamercolorbox}[wd=.4\paperwidth,ht=2.5ex,dp=1.125ex,leftskip=.3cm,rightskip=.3cm]{author in head/foot}%
  \veranstaltung
  \end{beamercolorbox}%
  \begin{beamercolorbox}[wd=.6\paperwidth,ht=2.5ex,dp=1.125ex,leftskip=.3cm,rightskip=.3cm plus1fil]{author in head/foot}%
    \usebeamerfont{author in head/foot}\semester\hspace*{5.2cm}\insertpagenumber
  \end{beamercolorbox}}%
  \vskip0pt%
  \footaline
}
\makeatletter


%%% time-saving commands

\newcommand{\okay}{\ding{51}}

\newcommand{\bilds}[1]{
	\begin{figure}%
	\centering%
	\fbox{\includegraphics[width=\textwidth]{#1}}
\end{figure}}

\newcommand{\bild}[3]{
	\begin{figure}%
	\centering%
	\includegraphics[width=#1\textwidth]{#2}
	\caption{#3}
\end{figure}}

\newcommand{\bildf}[3]{
	\begin{figure}%
	\centering%
	\fbox{\includegraphics[width=#1\textwidth]{#2}}
	\caption{#3}
\end{figure}}

\newcommand{\bildfc}[2]{
	\begin{figure}%
	\centering%
	\fbox{\includegraphics[width=#1\textwidth]{#2}}
\end{figure}}


\usepackage{attachfile}
\newcommand{\ta}[1]{\textattachfile[color=1 0 0]{#1}{\faFilePowerpointO}}

\newcommand{\pypy}[2]{\lstinputlisting[language={Python},caption={#1 \ta{#2}}]{#2}}


\hypersetup{
  pdftitle={<Lecture Title>},
  pdfauthor={<Lecturer>},
  pdfsubject={<Subject>},
  pdfkeywords={<Keywords>}
}

\author{<Lecturer>}
\title{<Lecture Title>}
\subtitle{\semester}
\date{\today}

